\section{Dissertation Organization}

\label{sec:organization} This dissertation is organized into the
following chapters.

\textbf{Chapter~\ref{chap:background}} provides technical background on
the device and architectural behaviors that determine write latency in
commercial ReRAM crossbar memories, including OxRAM bipolar switching
dynamics, program-verify loops, crossbar array architecture, and linear
block codes.

\textbf{Chapter~\ref{chap:brewrc}} presents BREW-RC, the first
non-invasive technique to recover the bit-exact ECC parity submatrix
from commercial ReRAM using write-access timing. It describes the
experimental characterization of write latency on two commercial
products, the SAT-based formulation of ECC recovery, and a validation of
the recovered parity matrix against destructive reverse engineering
results.

\textbf{Chapter~\ref{chap:trng}} applies the structural understanding
developed in Chapter~\ref{chap:brewrc} to the characterization of
stochastic write latency in commercial ReRAM. It demonstrates that prior
entropy claims are overstated due to unaccounted ECC activity, presents
an ECC-aware framework for entropy source construction and evaluation,
characterizes the effects of device aging on entropy yield, and
demonstrates a TRNG achieving 17$\times$ throughput improvement over
prior methods.

% \textbf{Chapter~\ref{chap:remanence}} connects the timing channels
% characterized in this dissertation to the broader class of hardware
% security threats that exploit physical device state. It reviews FPGA
% Pentimenti as the most developed instance of this class, identifies the
% open research challenges specific to ReRAM data remanence, and
% discusses what the contributions of this dissertation enable and what
% remains to be demonstrated.

\textbf{Chapter~\ref{chap:conclusion}} summarizes the contributions of
this dissertation, reflects on the implications for the security of
deployed ReRAM products, and identifies directions for future work.